\subsection{An exact ordinary noncapitulation and the twenty-seven S-corrections}

\begin{statement}{E.81}{Proved}{The prime-seventy-three norm polynomial and a finite S-capitulation reduction}
Put $F=B_1$, $E=B_2$ in E.76--E.79, with
$\beta_2^3=\beta_1$, and let
$G=\langle\sigma\rangle=\operatorname{Gal}(E/F)$.
Order the primes of $F$ above three and seventeen as in E.79.
Write $Q_1,Q_2,Q_3,T_1,T_2,T_3$ for the respective unique
primes of $E$ above them. Let $\mathfrak p$ be the actual
degree-one prime above $73$ with $t=a-\xi\mapsto33$ and
$\beta_1\mapsto22$. Its certified class in
$\operatorname{Cl}(F)=(\mathbb Z/3)^2$ is $(2,1)$.
Set $I=\mathfrak p\mathcal O_E$.

Use the actual genus Artin map from E.76,
\[
 \varphi:\operatorname{Cl}(E)\{3\}\longrightarrow
 \Gamma=\mathbb Z/9\oplus(\mathbb Z/3)^3,
 \tag{E81.1}
\]
where $\Gamma$ is the Galois group over $F$ of the maximal
abelian subfield of the three-primary Hilbert class field of
$E$. The image of $\varphi$ is the subgroup fixing $E$.
In the specified coordinates its values on the six upper
ramified primes and on $I$ are
\[
 C=\begin{pmatrix}
 3&3&3&0&3&6\\
 1&1&1&2&2&2\\
 0&0&0&2&2&2\\
 1&1&1&0&0&0
 \end{pmatrix},\qquad
 \varphi[I]=(6,0,0,0).
 \tag{E81.2}
\]
The rows are reduced modulo $9,3,3,3$, respectively.
Consequently $[I]$ has order exactly three in the ordinary
class group. More generally, every class of $F$ whose second
coordinate is nonzero retains order three after extension to
$E$; the ordinary capitulation kernel is contained in the
first coordinate subgroup.

There is an explicit comparison to the S-class question in
E.79. Choose the prime $\mathfrak p_2$ above $\mathfrak p$
with $\beta_2\mapsto63$ in $\mathbb F_{73}$. Then
\[
 [I]=(1+\sigma+\sigma^2)[\mathfrak p_2]
   \quad\hbox{in }\operatorname{Cl}(E),\qquad
 \pi_S[I]=j_S\pi_S[\mathfrak p],
 \tag{E81.3}
\]
where $\pi_S$ is the quotient by the classes of primes above
$3$ and $17$, and $j_S$ is actual ideal extension on S-classes.
The first element in (E81.3) is nonzero. Its image under the
second displayed quotient is evaluated by the following
finite set of ordinary principal-ideal tests.

For $a,b,c\in\{-1,0,1\}$, define the actual fractional ideals
\[
 J_{a,b,c}=
 I Q_1^{3a+2b+2c}Q_2^bQ_3^cT_1^{-1}T_2.
 \tag{E81.4}
\]
Then $j_S\pi_S[\mathfrak p]=0$ precisely when at least one of
these twenty-seven ideals is principal. Every $J_{a,b,c}$ is
$G$-invariant, has a principal cube, and has zero image under
$\varphi$. The certificate retains an actual $f_{a,b,c}\in F$
with
\[
 (f_{a,b,c})=
 \mathfrak p^3\prod_{i=1}^6\mathfrak s_i^{e_i(a,b,c)},
 \qquad J_{a,b,c}^{3}=f_{a,b,c}\mathcal O_E,
 \quad e(a,b,c)=(3a+2b+2c,b,c,-1,1,0).
 \tag{E81.5}
\]
Here $\mathfrak s_i$ are the six base primes in the prescribed
order. Both ideal identities are checked in the actual maximal
orders, including negative exponents.

The following lattices and literal principal functions make
the reduction to (E81.4) exact. With the valuation matrix $V$
of E.79, the seven proved S-principal relations in $E$ are
the columns of
\[
 M=\left(3V\ \middle|\ (0,0,0,1,1,1)^{\mathsf t}\right),
 \qquad
 H=\begin{pmatrix}
 9&6&6&3&3&0\\0&3&0&0&0&0\\0&0&3&0&0&0\\
 0&0&0&3&0&1\\0&0&0&0&3&1\\0&0&0&0&0&1
 \end{pmatrix},
 \quad H\mathbb Z^6=M\mathbb Z^7.
 \tag{E81.6}
\]
The six columns of $H$ are the divisors of the included
$s_1,\ldots,s_5$ from E.79 and of $\beta_2$, respectively.
Thus an element $Hq$ is the actual principal divisor of
$s_1^{q_1}\cdots s_5^{q_5}\beta_2^{q_6}$.
The quotient by these proved relations is
$\mathbb Z/9\oplus(\mathbb Z/3)^4$, of order $729$.
No assertion that these are all S-principal relations is
needed.

For an arbitrary correction exponent $e\in\mathbb Z^6$,
the genus obstruction vanishes exactly on the affine lattice
\[
 \begin{split}
 e&\in p+K\mathbb Z^6,\qquad p=(0,0,0,-1,1,0)^{\mathsf t},\\
 K&=\begin{pmatrix}
 3&2&2&0&0&0\\0&1&0&0&0&0\\0&0&1&0&0&0\\
 0&0&0&3&0&1\\0&0&0&0&3&1\\0&0&0&0&0&1
 \end{pmatrix},\qquad
 K\mathbb Z^6/H\mathbb Z^6\simeq(\mathbb Z/3)^3.
 \end{split}
 \tag{E81.7}
\]
To recover its representative (E81.4), compute
$z=K^{-1}(e-p)$ and take the centered residues of
$(z_1-z_4-z_5,z_2,z_3)$ modulo three as $(a,b,c)$.
The vector $q=H^{-1}(e-e(a,b,c))$ is integral, and the
principal function following (E81.6) is the explicit
comparison of the two ideals. All these matrices, their
Smith transformations and their inverses are retained.

There is a further numerical bound on these remaining
tests. Put
\[
 W=\ker\!\left(
    \operatorname{Cl}(E)\{3\}^{G}
       \xrightarrow{\ \varphi\ }\Gamma\right).
 \tag{E81.8}
\]
Then $|W|$ divides three, and all twenty-seven classes in
(E81.4) lie in $W$. Their class map is an affine map from
$(\mathbb Z/3)^3$ to this actual group. Accordingly the
number of principal ideals in the displayed list is one
of $0,9,27$.
The subgroup generated by the six S-prime classes has
order $27$ or $81$. Equivalently, the kernel of the actual
map $\mathbb Z^6\to\operatorname{Cl}(E)$,
$e\mapsto[\prod_i\mathfrak S_i^{e_i}]$, contains
$H\mathbb Z^6$ with index $27$ or $9$, respectively.

The remaining S-capitulation value, the higher annihilator
ideals of E.79, and the eventual exponent $h$ of E.69 are
not evaluated here. The bounded reduction searches retained
with the certificate found no positive generator; no
nonprincipality conclusion is drawn from those searches.
The ordinary noncapitulation above follows from (E81.2).
The exact positional comparisons remain E.46, E.72 and
E76.11. The manuscript anchor is LM~\S51.16,
fols.~733v--736v, scans~172--178.
\end{statement}

\begin{proof}
We describe the actual number-field maps first. The
certified base has primitive element $t=a-\xi$,
$a^3=17$, $\xi^3=\zeta$, and $\beta_1=a\xi$.
The ideal $(73,t-33)$ has norm $73$, and reduction gives
$\xi=55$, $a=15$, $\beta_1=22$ in its residue field.
Its class is $(2,1)$ in the certified class basis.
The program retains $f\in F$ and verifies
$(f)=\mathfrak p^3$ and $N_{F/\mathbb Q}f=73^3$.
It constructs $E=F[Y]/(Y^3-\beta_1)$ by an actual
relative maximal order. The absolute degree is $54$,
the discriminant is $-3^{135}17^{48}$, and the relative
power-basis index is one. At every base prime in $S$
it checks
\[
 \mathfrak s_i\mathcal O_E=\mathfrak S_i^{3},
 \qquad
 (e(\mathfrak S_i/\mathbb Q),f(\mathfrak S_i/\mathbb Q))
 =\begin{cases}(18,1)&i\le3,\\(9,2)&i\ge4.
 \end{cases}
 \tag{E81.9}
\]
These identities fix the prime ordering, rather than
matching two unordered lists. No class-group computation
in degree $54$ is used. The official maximal-order,
ideal-extension and principal-generator semantics are
those in \cite{E_PARICertifiedUnits}; the complete base
BNF is certified with default flag zero and return
value one.

Let $\mathcal G/F$ be the genus field defining $\Gamma$
in E.76. It is unramified over $E$, so Frobenius gives
the homomorphism (E81.1) on every ordinary ideal class
of $E$. At a prime outside $S$, restriction of
Frobenius to $\mathcal G/F$ is the Artin symbol of the
ideal norm down to $F$. In particular, $\mathfrak p$
splits in $E$: the root $63$ satisfies $63^3=22$ in
$\mathbb F_{73}$, and $\mu_3$ is contained there.
Thus $I$ is the product of its three conjugate upper
primes. This proves (E81.3) and
\[
             \varphi[I]=3\operatorname{Art}_{\mathcal G/F}
                                      (\mathfrak p).
 \tag{E81.10}
\]

The local normalization at the six ramified primes is
also explicit. At a three-prime use the upper uniformizer
$\beta_2+1$, whose norm is $\pi=\beta_1+1$; at a
seventeen-prime use $\beta_2$, whose norm is $\beta_1$.
These are exact positive norm identities because the
relative degree is odd. Arithmetic local reciprocity
and its norm/inclusion square therefore compute the
upper Frobenius as the base local Artin symbol of that
specified norm \cite[III Theorems~3.1 and~3.4]{E_MilneCFT}.

For clarity, these local symbols are represented by
actual coprime ideals in the certificate. Take the
modulus with exponents $(10,10,10,1,1,1)$ of E.76.
For its $i$th prime, choose by the Chinese remainder
map an element $x_i$ congruent to the specified
uniformizer one precision step beyond that modulus,
and congruent to one at its other primes. Then
$(x_i)\mathfrak s_i^{-1}$ is coprime to $S$, and the
desired local symbol is the negative of its ray Artin
class. The full integer quotient map of E.76 is then
applied to that class. This gives exactly the six
columns in (E81.2). Applying it to (E81.10) gives
$(6,0,0,0)$. All integer Smith maps and inverse maps
are checked again in the new certificate.

The ideal $I^3$ is generated by the included $f$, so
its class has order dividing three. Its nonzero Artin
image has order three, proving the claimed exact
order. The first base class generator is a prime
above three, with class $(1,0)$. Its extension is
$Q_1^3$, whose image is three times the first column
of (E81.2), hence zero. Since $(2,1)$ has image
$(6,0,0,0)$, extension followed by $\varphi$ on the
whole base class group is
\[
 (a,b)\longmapsto(6b,0,0,0).
 \tag{E81.11}
\]
This proves the assertion about all six classes with
nonzero second coordinate.

We next prove that the finite correction list is
complete through actual principal functions. Include
the six complete base S-unit generators of E.79 into
$E$. Formula (E81.9) multiplies each valuation vector
by three, giving $3V$. The divisor of $\beta_2$ is
$T_1+T_2+T_3$, giving the last column of $M$.
The columns of $H$ in (E81.6) are columns
$1,2,3,4,5,7$ of $M$. The omitted sixth column equals
the first minus the fourth and fifth plus three
times the seventh. This proves equality of the
two relation lattices and supplies the explicit
function for each column, in both directions.
The complete Smith transformation gives
$\mathbb Z^6/H\mathbb Z^6=
\mathbb Z/9\oplus(\mathbb Z/3)^4$.

For a correction $e$, the Artin image of
$I\prod\mathfrak S_i^{e_i}$ is
$(6,0,0,0)+Ce$. Reducing the four rows of (E81.2)
gives exactly
\[
 e_1+e_2+e_3\equiv0,\qquad
 e_4-e_6\equiv-1,\qquad
 e_5-e_6\equiv1\pmod3.
 \tag{E81.12}
\]
The solution lattice is $p+K\mathbb Z^6$. Direct
integer multiplication gives
\[
 K^{-1}H=
 \begin{pmatrix}
 3&0&0&1&1&0\\0&3&0&0&0&0\\0&0&3&0&0&0\\
 0&0&0&1&0&0\\0&0&0&0&1&0\\0&0&0&0&0&1
 \end{pmatrix}.
 \tag{E81.13}
\]
This proves the quotient and the inverse coordinate
rule in (E81.7). It gives exactly (E81.4) as the
twenty-seven representatives. If an S-principalization
is given with any integer exponents, its Artin image
vanishes, so (E81.12) applies. Reducing those exponents
by (E81.13) and multiplying by the explicit principal
function associated to $Hq$ gives a principal ideal
in (E81.4). Conversely any principal ideal in that
list exhibits $[I]$ as a combination of S-prime
classes, so its S-class is zero. This proves the
two directions with the actual ideal comparison.

For the cube identities, the six base prime classes
are $(1,0)$ three times and $(2,0)$ three times.
For $e=e(a,b,c)$ their weighted sum is zero modulo
three. Consequently
$\mathfrak p^3\prod\mathfrak s_i^{e_i}$ is principal
in $F$. Its certified principal generator is
$f_{a,b,c}$ in (E81.5). Cubing (E81.4) and using
(E81.9) proves the upper identity. The program
checks that identity in the upper maximal order,
while an independent rational-matrix computation
checks every lower principal-ideal identity in the
actual integral basis of $F$. The twenty-seven
exponent classes are pairwise distinct modulo $H$.

Finally E.76 computes the entire ordinary ambiguous
class group order as
\[
 |\operatorname{Cl}(E)\{3\}^{G}|
       =\frac{9\cdot3^6}{3\cdot27}=81.
 \tag{E81.14}
\]
This uses the certified base class number nine, the
six ramification indices three, and the ordinary-unit
field-norm index twenty-seven proved there. All six
ramified prime classes are $G$-invariant. Their
Artin columns span a subgroup of order twenty-seven,
as is verified both by the retained Smith maps and
by the three independent columns consisting of the
first, fourth, and fifth minus fourth columns.
Therefore the kernel (E81.8) has order dividing
three. The subgroup generated by the six S-prime
classes has Artin image of order twenty-seven and
kernel a subgroup of $W$, so its order is twenty-seven
or eighty-one. The kernel of its actual divisor-class
map has that index in $\mathbb Z^6$; comparison with
$[\mathbb Z^6:H\mathbb Z^6]=729$ proves the stated
index of additional principal S-relations.
Each candidate ideal is $G$-invariant and
satisfies (E81.12), so its class lies in $W$.
The principal functions for $H$ identify differences
of candidate classes with a homomorphism
$(\mathbb Z/3)^3\to W$. Translation by $[J_{0,0,0}]$
gives their affine class map. Such a map has either
zero, nine, or twenty-seven zero values, including
the case $W=0$. This proves the final bound.

For reproducibility, the degree-fifty-four computation
constructs only a maximal order and exact ideals.
The complete class-and-unit computation is the
already certified degree-eighteen base. The separate
unfiltered search was stopped after the genus
constraint was evaluated; its real interrupted
output is retained. The subsequent $108$ reductions
tested all twenty-seven necessary correction classes
in four directions each and found no generator.
Those outcomes are recorded as attempts. They do
not evaluate the zero set of the affine class map
just proved, and do not change the remaining
S-capitulation or universal-norm index in E.79.
\end{proof}
